
\chapter*{Abstract}
\addcontentsline{toc}{chapter}{Abstract}
\thispagestyle{plain}

Convolutional networks reach near-perfect accuracy on curated plant disease
datasets, and that number is routinely reported as if it described the field. This
project asks what survives when the photographs stop being curated.

We train four architectures under one frozen protocol on PlantVillage, measure the
same checkpoints on PlantDoc, and then test explanations for the difference instead
of asserting one. On 10,709 held-out curated images the best model reaches a
macro-F1 of 0.9931. On 236 field images covering 27 scored classes the same model
reaches 0.2232. The drop is 76.8 points, and it replicates across seeds.

We then asked what the decision depends on, by intervention rather than by
inspecting saliency maps. Replacing an image's leaf with a leaf from another class
moved the prediction to the donor class 50.3\,\% of the time; replacing the
background moved it 21.8\,\%, against a 2.6\,\% chance rate and a 4.4\,\% control
for the edit itself. The leaf dominates, so the model has not learned to ignore the
plant, but background content can still steer it one time in five.

That measurement predicts the outcome of the fix we tested next. Background
augmentation during training moved field macro-F1 by 1.24 points, with a 95\,\%
confidence interval from $-3.32$ to $+5.80$ and one seed in three going backwards,
while costing 0.87 points in the lab. Recovering 1.6\,\% of the gap is not a
detectable effect.

The gap is therefore not background overfitting that augmentation can repair. It is
leaf appearance that differs between curated and field photographs. Confidence
thresholds were also measured and do not repair it: abstention reduces coverage
without producing an operating point an advisor could act on.

The contribution is a measured, reproducible statement of what a clean-image
classifier is worth in the field, with the mechanism tested rather than assumed,
and with every reported number traceable to a stored artifact.

\vspace{0.6cm}
\noindent\textbf{Keywords:} plant disease detection, convolutional neural networks,
domain shift, field validation, Grad-CAM, calibration, background augmentation
